% ============================================================
%  SuVa Protocol — Phase 2 : ETHFalcon on EVM
%  Rapport technique
%  Auteur : Mohamed Lamine MBENGUE
%  Portfolio : https://portfolio-alamine.web.app
%  Date : Mars 2026
% ============================================================
\documentclass[11pt,a4paper]{article}

% ── Encodage & langue ────────────────────────────────────────────────────────
\usepackage[utf8]{inputenc}
\usepackage[T1]{fontenc}
\usepackage[french]{babel}

% ── Typographie ──────────────────────────────────────────────────────────────
\usepackage{lmodern}
\usepackage{microtype}
\usepackage{fourier}          % Utopia font — cohérent avec les slides

% ── Mise en page ─────────────────────────────────────────────────────────────
\usepackage[top=2.5cm,bottom=2.5cm,left=2.8cm,right=2.8cm]{geometry}
\usepackage{parskip}
\setlength{\parindent}{0pt}

% ── Maths ────────────────────────────────────────────────────────────────────
\usepackage{amsmath,amssymb,amsthm}
\usepackage{physics}

% ── Couleurs & boîtes ────────────────────────────────────────────────────────
\usepackage[dvipsnames,table]{xcolor}

\definecolor{suvaBlue}{RGB}{0,82,155}
\definecolor{suvaGold}{RGB}{204,153,0}
\definecolor{suvaGreen}{RGB}{0,128,80}
\definecolor{suvaRed}{RGB}{180,30,30}
\definecolor{boxBg}{RGB}{240,248,255}
\definecolor{warnBg}{RGB}{255,248,230}

\usepackage{tcolorbox}
\tcbuselibrary{breakable,skins}

% Boite résultat (vert/bleu)
\newtcolorbox{resultbox}[1][]{
  enhanced,
  breakable,
  colback=boxBg,
  colframe=suvaBlue,
  fonttitle=\bfseries\color{white},
  coltitle=white,
  attach boxed title to top left={yshift=-2mm,xshift=4mm},
  boxed title style={colback=suvaBlue},
  title={Résultat},
  #1
}

% Boite avertissement (orange)
\newtcolorbox{warningbox}[1][]{
  enhanced,
  breakable,
  colback=warnBg,
  colframe=suvaGold,
  fonttitle=\bfseries\color{white},
  coltitle=white,
  attach boxed title to top left={yshift=-2mm,xshift=4mm},
  boxed title style={colback=suvaGold},
  title={Important},
  #1
}

% ── Code source ──────────────────────────────────────────────────────────────
\usepackage{listings}
\lstset{
  basicstyle=\ttfamily\small,
  keywordstyle=\color{suvaBlue}\bfseries,
  commentstyle=\color{gray},
  stringstyle=\color{suvaGreen},
  numbers=left,
  numberstyle=\tiny\color{gray},
  breaklines=true,
  frame=single,
  rulecolor=\color{suvaBlue!40},
  backgroundcolor=\color{suvaBlue!5},
  captionpos=b
}

% ── Tableaux ─────────────────────────────────────────────────────────────────
\usepackage{booktabs,multirow,array}

% ── Liens ────────────────────────────────────────────────────────────────────
\usepackage{hyperref}
\hypersetup{
  colorlinks=true,
  linkcolor=suvaBlue,
  urlcolor=suvaBlue,
  citecolor=suvaGreen
}

% ── En-têtes et pieds de page ────────────────────────────────────────────────
\usepackage{fancyhdr}
\pagestyle{fancy}
\fancyhf{}
\fancyhead[L]{\small\color{suvaBlue}\textbf{SuVa Protocol — Phase 2}}
\fancyhead[R]{\small\color{gray}ETHFalcon on EVM}
\fancyfoot[C]{\thepage}
\renewcommand{\headrulewidth}{0.4pt}
\renewcommand{\footrulewidth}{0pt}

% ── Sections colorées ────────────────────────────────────────────────────────
\usepackage{titlesec}
\titleformat{\section}
  {\Large\bfseries\color{suvaBlue}}
  {\thesection.}{0.6em}{}[\color{suvaBlue}\titlerule]
\titleformat{\subsection}
  {\large\bfseries\color{suvaBlue!80}}
  {\thesubsection.}{0.5em}{}

% ── Abréviations ─────────────────────────────────────────────────────────────
\newcommand{\ethfalcon}{\textsc{ETHFalcon}}
\newcommand{\ethdilithium}{\textsc{ETHDilithium}}
\newcommand{\Zq}{\mathbb{Z}_{q}}
\newcommand{\ring}{\mathbb{Z}_q[X]/(X^n+1)}

% ============================================================
%  DOCUMENT
% ============================================================
\begin{document}

% ── Page de titre ────────────────────────────────────────────────────────────
\begin{titlepage}
  \centering
  \vspace*{1.5cm}

  {\color{suvaBlue}\rule{\linewidth}{2pt}}
  \vspace{0.5cm}

  {\Huge\bfseries\color{suvaBlue} SuVa Protocol}\\[0.3cm]
  {\LARGE\color{suvaGold} Rapport Technique — Phase 2}\\[0.2cm]
  {\Large ETHFalcon : Falcon sur EVM avec KeccakPRNG}

  \vspace{0.5cm}
  {\color{suvaBlue}\rule{\linewidth}{2pt}}

  \vspace{2cm}

  \begin{tabular}{rl}
    \textbf{Auteur :}    & Mohamed Lamine MBENGUE \\[4pt]
    \textbf{Portfolio :} & \href{https://portfolio-alamine.web.app}{portfolio-alamine.web.app} \\[4pt]
    \textbf{Projet :}    & SuVa Protocol — Post-Quantum Signatures on EVM \\[4pt]
    \textbf{Phase :}     & Phase 2 — ETHFalcon \\[4pt]
    \textbf{Date :}      & Mars 2026 \\[4pt]
  \end{tabular}

  \vfill

  \begin{tcolorbox}[colback=suvaBlue!8,colframe=suvaBlue,width=0.85\linewidth]
    \centering
    \textbf{Résumé exécutif}\\[6pt]
    La Phase 2 du projet SuVa Protocol implémente \ethfalcon{}, une variante
    de la signature post-quantique Falcon où SHAKE-256 est remplacé par
    KeccakPRNG (Keccak-256). Cette modification rend la vérification
    entièrement compatible avec l'EVM. Le contrat \texttt{ZKNOX\_ethfalcon.sol}
    effectue la vérification en \textbf{$\approx 46{,}67$ M gas} (multiplication
    schoolbook $O(n^2)$, $n=256$), contre 6,58 M pour \ethdilithium{} et
    13,42 M pour Dilithium standard.
  \end{tcolorbox}

  \vspace{1cm}
  {\color{suvaBlue}\rule{\linewidth}{1pt}}
\end{titlepage}

\tableofcontents
\newpage

% ============================================================
\section{Résumé exécutif}
% ============================================================

\begin{resultbox}[title={Résultats de la Phase 2}]
\begin{center}
\begin{tabular}{lrrr}
\toprule
\textbf{Schéma} & \textbf{Gas (verify)} & \textbf{Taille sig.} & \textbf{Taille vk} \\
\midrule
Dilithium (NIST II) & 13{,}42 M & $\approx 2\,420$ B & $\approx 1\,312$ B \\
ETHDilithium (Keccak) & 6{,}58 M & $\approx 2\,420$ B & $\approx 1\,312$ B \\
\textbf{ETHFalcon-256} & \textbf{46{,}67 M} & \textbf{356 B} & \textbf{448 B} \\
\bottomrule
\end{tabular}
\end{center}
\end{resultbox}

\medskip

La Phase 2 atteint ses objectifs principaux :
\begin{itemize}
  \item Le module Python \texttt{ETHFalcon256} (keygen / sign / verify) fonctionne
        correctement et produit des vecteurs de test reproductibles.
  \item Le contrat Solidity \texttt{ZKNOX\_ethfalcon.sol} vérifie correctement les
        signatures ETHFalcon-256, validé par deux tests Forge (message correct et
        message erroné).
  \item La taille de signature ETHFalcon (356 B) est \textbf{6,8× plus petite}
        que Dilithium ($\approx 2420$ B), réduisant le calldata.
  \item Le coût en gas plus élevé ($\approx 46{,}67$ M vs $6{,}58$ M) est dû à la
        multiplication polynomiale schoolbook $O(n^2)$ ; une NTT dédiée à
        $q=12289$ réduirait ce coût (objectif Phase 3).
\end{itemize}

% ============================================================
\section{Contexte cryptographique — Falcon}
% ============================================================

\subsection{Principe général}

Falcon est un schéma de signature basé sur les réseaux NTRU, standardisé par
le NIST (FIPS 206, 2024) \cite{NIST_PQC}. Il repose sur le problème difficile
\textbf{NTRU-SIS} : étant donné une paire $(f,g)$ de petits polynômes vérifiant
$fh \equiv g \pmod{q}$, retrouver $f$ et $g$.

\textbf{Paramètres Falcon-256 :}
\begin{center}
\begin{tabular}{ll}
\toprule
Paramètre & Valeur \\
\midrule
$n$ (degré du polynôme)        & 256 \\
$q$ (module)                   & 12289 \\
Anneau                         & $\ring$ \\
Borne de norme $\beta^2$       & 16\,468\,416 \\
Longueur sel (\texttt{SALT\_LEN}) & 40 octets \\
Taille signature totale        & 356 octets \\
Taille clé publique            & 448 octets \\
\bottomrule
\end{tabular}
\end{center}

\subsection{Schéma de vérification}

La vérification d'une signature Falcon utilise l'équation :
\[
  \|s_0\|^2 + \|s_1\|^2 \leq \beta^2
\]
où les vecteurs sont calculés comme suit :
\begin{enumerate}
  \item $t \leftarrow \mathrm{HashToPoint}(\mathrm{salt} \| \mathrm{msg})$ : un
        polynôme de $n=256$ coefficients dans $\Zq$.
  \item $s_0 \leftarrow t - h \cdot s_1 \pmod{q,\, X^n+1}$ : erreur de signature.
  \item Vérifier $\|s_0\|^2 + \|s_1\|^2 \leq \beta^2$ avec centrage modular.
\end{enumerate}

Les normes sont calculées avec les coefficients \textbf{centrés} dans
$(-q/2, q/2]$ : si $c > q/2$, on traite $c$ comme $c - q$.

\subsection{Comparaison avec Dilithium}

\begin{center}
\begin{tabular}{lcc}
\toprule
Propriété & Dilithium (NIST II) & Falcon-256 \\
\midrule
Problème dur       & MLWE + MSIS & NTRU-SIS \\
Taille de signature & $\approx 2420$ B & $356$ B \\
Clé publique        & $\approx 1312$ B & $448$ B \\
HashToPoint         & SHAKE-256 & SHAKE-256 (standard) \\
\bottomrule
\end{tabular}
\end{center}

% ============================================================
\section{Modification ETHFalcon — Keccak à la place de SHAKE-256}
% ============================================================

\subsection{Motivation}

L'opcode \texttt{KECCAK256} (0x20) est natif sur l'EVM avec un coût de 30 gas
de base + 6 gas par mot de 32 octets. En revanche, SHAKE-256 n'est pas
disponible nativement et nécessite une implémentation coûteuse en bytecode
Solidity.

La fonction \texttt{HashToPoint} est le seul endroit où Falcon utilise un XOF
(eXtendable Output Function). Remplacer SHAKE-256 par un KeccakPRNG (déjà
utilisé dans ETHDilithium) rend la vérification entièrement compatible EVM.

\subsection{Algorithme HashToPoint avec KeccakPRNG}

\begin{warningbox}[title={Compatibilité Python / Solidity}]
L'algorithme ci-dessous doit être identique en Python (module de référence) et
en Solidity (contrat on-chain). Toute divergence produirait une vérification
incorrecte.
\end{warningbox}

\textbf{Protocole KeccakPRNG utilisé :}
\begin{enumerate}
  \item Calculer la graine : $\texttt{seed} = \mathrm{keccak256}(\mathrm{salt} \| \mathrm{msg})$
  \item Boucle de génération (compteur $c$ initialisé à 0) :
    \begin{enumerate}
      \item Générer un bloc de 32 octets :
            $\texttt{block} = \mathrm{keccak256}(\texttt{seed} \| c_{\mathrm{BE64}})$
      \item Traiter le bloc par paires d'octets $(b_0, b_1)$ en \textbf{big-endian} :
            $\mathrm{elt} = (b_0 \ll 8) \mid b_1$
      \item Rejection sampling : si $\mathrm{elt} < 5q = 61445$, alors
            $t[\mathrm{filled}] = \mathrm{elt} \bmod q$, $\mathrm{filled}$++
      \item Incrémenter $c$
    \end{enumerate}
  \item Continuer jusqu'à $\mathrm{filled} = n = 256$
\end{enumerate}

\textbf{Note sur le byte order :} Falcon standard utilise big-endian dans
\texttt{falcon\_sign.py} (ligne : \texttt{elt = (twobytes[0] << 8) + twobytes[1]}).
ETHFalcon conserve ce choix, ce qui diffère de certaines implémentations
little-endian.

% ============================================================
\section{Implémentation}
% ============================================================

\subsection{Structure des fichiers}

\begin{center}
\begin{tabular}{lp{8cm}}
\toprule
Fichier & Rôle \\
\midrule
\texttt{src/ZKNOX\_falcon\_params.sol} &
  Constantes : $q$, $n$, $\beta^2$, seuil rejection, longueurs \\
\texttt{src/ZKNOX\_falcon\_poly.sol} &
  Multiplication polynomiale schoolbook + norme L2 centrée \\
\texttt{src/ZKNOX\_ethfalcon.sol} &
  Contrat principal, fonction \texttt{verify()} \\
\midrule
\texttt{python-ref/\ldots/Falcon\_falcon/falcon\_eth.py} &
  Classe \texttt{ETHFalcon256} (Python de référence) \\
\texttt{python-ref/\ldots/generate\_falcon\_test\_vectors.py} &
  Générateur de vecteurs de test Solidity \\
\texttt{test/ZKNOX\_ethfalcon.t.sol} &
  Tests Forge (généré automatiquement) \\
\bottomrule
\end{tabular}
\end{center}

\subsection{Module Python}

Le module Python suit la même architecture qu'ETHDilithium :

\begin{lstlisting}[language=Python,caption={Utilisation d'ETHFalcon256}]
from Falcon_dilithium_py.Falcon_falcon.falcon_eth import ETHFalcon256

msg = b"SuVa Protocol Phase 2"
sk, vk = ETHFalcon256.keygen()
sig = ETHFalcon256.sign(sk, msg)
ok = ETHFalcon256.verify(vk, msg, sig)
assert ok  # True
\end{lstlisting}

Le module sous-classe \texttt{falcon\_sign.Falcon} en remplacant uniquement
la méthode \texttt{\_\_hash\_to\_point\_\_} par \texttt{hash\_to\_point\_keccak}.
Cela garantit que la signature (opération de sign) utilise aussi KeccakPRNG,
assurant la cohérence sign/verify.

\subsection{Contrat Solidity — \texttt{ZKNOX\_ethfalcon.sol}}

L'interface du contrat accepte les données pré-décodées (même approche
qu'\texttt{ZKNOX\_ethdilithium.sol}) :

\begin{lstlisting}[language=Solidity,caption={Interface verify()}]
function verify(
    uint256[256] calldata h,    // 256 coefficients cle publique (0..q-1)
    bytes        calldata salt, // 40 octets du sel de signature
    uint256[256] calldata s1,   // 256 coeff s1 decodesmodq (0..q-1)
    bytes        calldata msg   // message original
) external pure returns (bool)
\end{lstlisting}

\textbf{Algorithme en 4 étapes :}
\begin{enumerate}
  \item \textbf{HashToPoint :} $t \leftarrow \texttt{\_hashToPoint}(\mathrm{salt}, \mathrm{msg})$
        via KeccakPRNG (décrit ci-dessus).
  \item \textbf{Multiplication polynomiale :}
        $hs_1 \leftarrow \texttt{poly\_mul}(h, s_1)$
        dans $\ring$ avec la règle $X^{256} \equiv -1$.
  \item \textbf{Calcul de $s_0$ :}
        $s_0[i] \leftarrow (t[i] + q - hs_1[i]) \bmod q$.
  \item \textbf{Vérification de norme :}
        $\|s_0\|^2 + \|s_1\|^2 \leq \beta^2 = 16\,468\,416$.
\end{enumerate}

\subsection{Génération des vecteurs de test}

Le script \texttt{generate\_falcon\_test\_vectors.py} :
\begin{enumerate}
  \item Génère une paire de clés et signe un message de test avec
        \texttt{ETHFalcon256}.
  \item Décode la clé publique ($h$ via \texttt{deserialize\_to\_poly}) et
        la signature ($s_1$ via \texttt{decompress}).
  \item Convertit $s_1$ signé en non-signé : \texttt{s1\_uint[i] = s1[i] \% q}.
  \item Écrit \texttt{test/ZKNOX\_ethfalcon.t.sol} avec les tableaux
        initialisés élément par élément (Solidity n'accepte pas les littéraux
        de grands tableaux).
\end{enumerate}

% ============================================================
\section{Mesures de gas}
% ============================================================

\begin{resultbox}[title={Résultats forge test --gas-report}]
\begin{center}
\begin{tabular}{lrrr}
\toprule
\textbf{Contrat} & \textbf{Gas verify} & \textbf{Taille sig.} & \textbf{Multiplication} \\
\midrule
\texttt{ZKNOX\_dilithium}    & 13{,}42 M & $\approx 2420$ B & NTT ($O(n\log n)$) \\
\texttt{ZKNOX\_ethdilithium} & 6{,}58 M  & $\approx 2420$ B & NTT ($O(n\log n)$) \\
\textbf{\texttt{ZKNOX\_ethfalcon}} & \textbf{46{,}67 M} & \textbf{356 B} & Schoolbook ($O(n^2)$) \\
\bottomrule
\end{tabular}
\end{center}
\end{resultbox}

\subsection{Analyse du coût en gas}

Le coût élevé de \texttt{ZKNOX\_ethfalcon} ($\approx 46{,}67$ M gas) est
principalement dû à la multiplication polynomiale schoolbook $O(n^2)$ :
$256 \times 256 = 65\,536$ multiplications modulo $q=12289$.

\textbf{Décomposition estimée :}
\begin{itemize}
  \item Multiplication \texttt{poly\_mul} : $\approx 40$ M gas (65\,536 modmul)
  \item HashToPoint (KeccakPRNG, $\approx 16$--20 blocs keccak) :
        $\approx 5$--6 M gas
  \item Norme et vérification : $< 1$ M gas
\end{itemize}

\subsection{Comparaison gaz/octet de signature}

\begin{center}
\begin{tabular}{lrrl}
\toprule
Schéma & Gas & Sig (B) & Gas/octet \\
\midrule
Dilithium     & 13{,}42 M & 2420 & $5\,545$ \\
ETHDilithium  & 6{,}58 M  & 2420 & $2\,719$ \\
ETHFalcon     & 46{,}67 M & 356  & $131\,101$ \\
\bottomrule
\end{tabular}
\end{center}

\begin{warningbox}[title={Observation}]
Malgré une signature 6,8× plus courte, ETHFalcon consomme 7,1× plus de gas
qu'ETHDilithium. L'introduction d'une NTT pour $q=12289$ (Phase 3) est
indispensable pour rendre ETHFalcon compétitif.
\end{warningbox}

% ============================================================
\section{Analyse de sécurité}
% ============================================================

\subsection{Sécurité de la construction ETHFalcon}

Le remplacement de SHAKE-256 par KeccakPRNG affecte uniquement la fonction
\texttt{HashToPoint}. Les propriétés de sécurité de Falcon reposent sur :

\begin{enumerate}
  \item \textbf{Résistance NTRU-SIS :} Inchangée. Le problème dur reste
        retrouver $(f,g)$ à partir de $h = gf^{-1} \bmod q$. Le changement
        de hash n'affecte pas ce problème.

  \item \textbf{Indistingabilité de la fonction de hachage :} KeccakPRNG est
        une construction pseudo-aléatoire basée sur Keccak-256, qui est
        considérée comme une fonction de hachage sûre. L'output est
        computationnellement indistingable d'un XOF aléatoire pour des
        adversaires polynomiaux.

  \item \textbf{Uniformité du HashToPoint :} La rejection sampling garantit
        une distribution uniforme sur $\Zq$ pour chaque coefficient de $t$.
        Le seuil $5q = 61445$ assure une probabilité d'acceptation de
        $5q/2^{16} \approx 93.7\%$, identique à l'implémentation SHAKE.

  \item \textbf{Séparabilité sign/verify :} La modification est transparente
        pour le signer (la signature ne dépend de hash que via le sel, déjà
        incorporé dans le processus Falcon gaussien).
\end{enumerate}

\subsection{Niveau de sécurité}

\begin{center}
\begin{tabular}{lcc}
\toprule
Type d'attaque & Falcon-256 standard & ETHFalcon-256 \\
\midrule
Quantique (bits) & $\sim 103$ & $\sim 103$ \\
Classique (bits) & $\sim 103$ & $\sim 103$ \\
\bottomrule
\end{tabular}
\end{center}

La sécurité est identique à Falcon-256 standard. La seule différence est
computationnelle (on-chain), pas cryptographique.

% ============================================================
\section{Tests et validation}
% ============================================================

\subsection{Tests Python}

\begin{lstlisting}[language=bash,caption={Exécution des tests Python}]
cd python-ref/
py -3 -m Falcon_dilithium_py.generate_falcon_test_vectors
# Output:
# [*] Generating ETHFalcon256 test vectors
#     verify(correct) = True  OK
#     verify(wrong)   = False  (expected False)
# [+] Written: .../test/ZKNOX_ethfalcon.t.sol
\end{lstlisting}

\subsection{Tests Forge}

\begin{lstlisting}[language=bash,caption={Résultats forge test}]
forge test --match-contract ETHFalconTest -vv --gas-report
# Ran 2 tests for test/ZKNOX_ethfalcon.t.sol:ETHFalconTest
# [PASS] testVerify()         (gas: 46732996)
# [PASS] testVerifyWrongMsg() (gas: 46734869)
# Suite result: ok. 2 passed; 0 failed; 0 skipped
\end{lstlisting}

Les deux tests passent :
\begin{itemize}
  \item \texttt{testVerify()} : vérifie qu'une signature valide est acceptée.
  \item \texttt{testVerifyWrongMsg()} : vérifie qu'un message incorrect est
        rejeté (même $h$, même sel, même $s_1$, message modifié).
\end{itemize}

% ============================================================
\section{Perspectives — Phase 3}
% ============================================================

\subsection{Objectif principal : NTT pour q = 12289}

La Phase 3 doit implémenter une Number Theoretic Transform (NTT) adaptée
à $q = 12289$ pour remplacer la multiplication schoolbook. Cela permettrait
de réduire la complexité de $O(n^2)$ à $O(n \log n)$.

\textbf{Défi :} $q = 12289 = 3 \cdot 2^{12} + 1$ est premier et admet une
racine primitive $\omega = 3$ d'ordre $q-1 = 12288 = 2^{12} \cdot 3$. Une
NTT pour $n=256 = 2^8$ est donc réalisable avec une racine $2n$-ième de l'unité.

\textbf{Gain attendu :} Réduction de $\approx 40$ M gas (schoolbook) à
$\approx 2$--4 M gas (NTT), rendant ETHFalcon compétitif avec ETHDilithium.

\subsection{QuantumVault.sol}

L'objectif final est \texttt{QuantumVault.sol}, un smart contract de coffre
post-quantique combinant :
\begin{itemize}
  \item \textbf{ETHFalcon NTT} pour les signatures compactes (356 B),
  \item \textbf{ETHDilithium} comme fallback ou double vérification,
  \item Un mécanisme d'upgradabilité conforme ERC-2535 (Diamond proxy).
\end{itemize}

\begin{center}
\begin{tabular}{lrr}
\toprule
Phase & Objectif & Statut \\
\midrule
Phase 1 & ETHDilithium (NTT, Keccak)     & Terminée (6,58 M gas) \\
Phase 2 & ETHFalcon (schoolbook, Keccak)  & Terminée (46,67 M gas) \\
Phase 3 & ETHFalcon (NTT, q=12289)        & Planifiée ($\sim 4$ M gas) \\
Phase 4 & QuantumVault.sol                & Planifiée \\
\bottomrule
\end{tabular}
\end{center}

% ============================================================
%  Références
% ============================================================
\newpage
\begin{thebibliography}{9}

\bibitem{NIST_PQC}
  NIST,
  \textit{FIPS 206: Module-Lattice-Based Digital Signature Standard (ML-DSA)},
  FIPS 206, 2024.
  \url{https://csrc.nist.gov/pubs/fips/206/final}

\bibitem{Falcon_Spec}
  P.-A. Fouque, J. Hoffstein, P. Kirchner, V. Lyubashevsky, T. Pornin,
  T. Prest, T. Ricosset, G. Seiler, W. Whyte, Z. Zhang,
  \textit{Falcon: Fast-Fourier Lattice-based Compact Signatures over NTRU},
  Specification v1.2, 2020.
  \url{https://falcon-sign.info/falcon.pdf}

\bibitem{ZKNOX}
  R. Dubois, S. Masson,
  \textit{ZKNOX: Ethereum-Friendly Post-Quantum Signatures},
  GitHub, 2025.
  \url{https://github.com/ZKNox}

\bibitem{KeccakPRNG}
  Z. Zhang,
  \textit{Falcon-go: Keccak PRNG implementation},
  GitHub.
  \url{https://github.com/zhenfeizhang/falcon-go}

\end{thebibliography}

\end{document}
